%!TEX root = ../thesis.tex

\section{はじめに}

二足歩行ロボットは，人間の生活環境と親和性の高い移動手段として，災害対応，屋外作業支援，日常環境での移動支援など，幅広い応用が期待されている．これまでの研究により，平坦かつ剛体の床面を対象とした歩行制御技術は大きく発展し，安定した二足歩行の実現が可能となってきた．

一方で，実環境においては，地表は砂地や土壌，芝生，伸縮性のあるスポンジやクッションのような床などの軟弱地面が多く存在し，これらの環境では足底の沈み込みや床変形が生じる．このような床変形は，支持脚の高さ変化を引き起こし，ロボットの重心高さが瞬間的かつ不連続に変動する要因となる．結果として，重心位置や運動量の予測誤差が増大し，歩行の不安定化や転倒につながるという問題がある．

従来の二足歩行制御では，線形倒立振子モデル（Linear Inverted Pendulum Model: LIPM）を用いた手法が広く用いられてきた．LIPMは重心高さを一定と仮定することで運動方程式を線形化し，計算効率の高い歩行生成を可能にしている．しかし，床沈み込みにより重心高さが変動する環境では，この仮定が成立せず，LIPMに基づく歩行制御の安定性が大きく損なわれる．

近年，変動重心高さを扱うモデルとして可変重心高さモデルやVHIP（Variable Height Inverted Pendulum）などが提案されているが，これらの手法はモデルの複雑化や計算負荷の増加を伴い，既存のLIPMベースの歩行制御器へ容易に導入できない場合が多い．そのため，軟弱地面に適応しつつ，既存手法との親和性を維持した歩行制御アルゴリズムの確立が求められている．